\section{Limitations}
\label{sec:limitations}

\begin{itemize}
  \item \textbf{Small models, two families.} All results are for 1.5--1.7B instruction-tuned models. Larger models hold more of each task's knowledge, which may move tasks between regimes. For example, a larger model may know enough about people that Twin-2K becomes a redirect task.
  \item \textbf{Run counts are still modest.} Under the standard brief, Fable has:
  \begin{itemize}
    \item on Qwen: three runs each on SST-5, SST-5 inverted and facts, and one on Twin-2K;
    \item on SmolLM2: one run each on SST-5 inverted and facts.
  \end{itemize}
  Opus has one run per Qwen task, and the preserve condition has one run per compile task. Four further Fable runs stopped part-way when a usage limit was reached and are excluded. A second Twin-2K Fable run (\texttt{r4}) finished but broke the no-fitting rule, and is reported separately (Appendix~\ref{app:deviations}).
  \item \textbf{LoRA seeds and tuning.}
  \begin{itemize}
    \item \textbf{Seeds.} SST-5 and SST-5 inverted have three seeds each. Twin-2K, facts and SmolLM2 have one, because those runs take hours.
    \item \textbf{Tuning.} LoRA used one default configuration everywhere except facts, where we searched until it learned.
    \item \textbf{The plateau.} A learning-rate warmup or larger learning rate might help it escape the plateau on SST-5 inverted. We report the failure as a property of the default configuration, not of LoRA in general.
  \end{itemize}
  \item \textbf{Dev use by programmers.} Programmers measured candidates on dev, and some chose among candidates by dev accuracy, or built constants from unlabelled dev activations. The auditors quantified this where possible; on facts it inflated the preserve run's dev gain from about 4 to 6 points. Test numbers, which no programmer saw, are the reliable comparison.
  \item \textbf{Rule boundaries.} Three permissions in the brief are judgement calls:
  \begin{itemize}
    \item label-derived tables, if disclosed;
    \item closed-form solves on weights or activation statistics, such as pseudo-inverses;
    \item one-dimensional grid searches over a scalar.
  \end{itemize}
  Several of the best compile-regime edits rely on the first, so a stricter rule would lower programming's results there.
  \item \textbf{Infrastructure.}
  \begin{itemize}
    \item Two Opus runs ended early because Kaggle GPU sessions were unavailable for hours.
    \item Queued GPU jobs were sometimes cancelled and re-pushed. The ledgers show no pass was charged twice, but one run probably executed a job more than once without charging the duplicate passes.
    \item Some evaluations ran on P100 rather than T4 GPUs. With fp16 inference the effect on letter probabilities is negligible, but the results are not bit-identical.
  \end{itemize}
  \item \textbf{Protocol slips by the orchestrator.} One SST-5 replication (\texttt{r5}) was told another run's result mid-run. Its audit dated its final design to six hours before the message, so we count it as independent. Several remote job outputs were deleted while unsticking the queue; local copies were kept.
\end{itemize}

\section{Conclusion}
Weight-space programming promises a new way to change a model: read the weights, work out the computation, and write the fix directly. We asked when that works and what it costs, holding everything fixed except how much the model already knew, and repeated the key experiments on a second model family.

The answer has four parts.
\begin{itemize}
  \item \textbf{Redirect.} When the model has the knowledge and misreads it, programmers reliably find tiny fixes that match fine-tuning and fail less often.
  \item \textbf{Compile.} When the answer is a table, programmers compile it into the weights and can beat fine-tuning. The most accurate way to do so bypasses the network and leaves the model at chance on general knowledge. An instruction to preserve the model changes the edit the programmer builds, at a cost in accuracy.
  \item \textbf{Learn.} When the improvement has to be learned, fine-tuning wins, even with far fewer labels.
  \item \textbf{Combine.} A programmed edit makes fine-tuning faster and more reliable, and at the same training budget the combination beats both.
\end{itemize}

In our runs, weight-space programming acted as a compiler for rules and tables and a starting point for learning, not a replacement for it. Anyone using or evaluating it should state what must be preserved and measure what an edit does to the rest of the model, with measures a temperature change cannot fake.

\paragraph{Reproducibility.} Code, task builders, the programmer briefs, every programmer's log, edit and final report, all audit reports, the per-run cost accounting and all result files are in a private repository (\url{https://github.com/Puneet-Bajaj-IITM/wsp-vs-lora}), to be made public with the preprint.

\paragraph{Use of AI.} The weight-space programming runs and their audits were carried out by Claude models, as described in Section~\ref{sec:design}. Claude was also used to build the experiment infrastructure and to draft this manuscript. The author reviewed all results and text.
